先说明前面误差界里各项的含义。$p_0$ 是数据分布，$p_T$ 是人为加噪后的终端分布，
$\pi$ 是实际开始生成时采用的先验。即使分数网络已经足够准确，
若 $p_T$ 与 $\pi$ 差别很大，从 $\pi$ 出发仍会留下误差。
对反向 SDE，在适当正则性、可积性和路径绝对连续条件下，
常见的控制形式是
\[
D_{\rm KL}(p_0\|p_0^\theta)
\leq D_{\rm KL}(p_T\|\pi)
 +\frac12\int_0^Tg(t)^2\mathbb E_{p_t}
 \|s_\theta(X_t,t)-\nabla\log p_t(X_t)\|^2\,\mathrm dt.
\]
它来自路径分布的相对熵控制，再利用取端点边缘不会增加 KL。
因此原稿前面出现的 $\lambda(t)$ 在这个结论中要取相应的似然权重，
不能任意取一个正权重就声称仍有同一个界。这里也暂未计入有限步采样的数值误差。
\cite{sde,theorybook}

对 ODE 的轨迹误差，可以从更直接的估计入手。
设真实速度 $v$ 关于空间是 $L$-Lipschitz，近似速度沿近似轨迹的误差不超过 $\delta(t)$。
若两条轨迹的距离为 $e(t)$，则积分形式给出
\[
e(t)\leq e(0)+\int_0^t\bigl(Le(r)+\delta(r)\bigr)\,\mathrm dr
\leq e^{Lt}e(0)+\int_0^te^{L(t-r)}\delta(r)\,\mathrm dr.
\]
这一步说明了局部预测误差如何随积分过程累积。
但训练损失给出的通常是分布平均下的均方误差，
不能不加条件就用它替换这里的逐轨迹界；轨迹距离也不自动等于密度的 KL 误差。
后面的采样加速是在这样的误差基础上减少数值求解成本，
并不能把模型误差一并消除。靠近零噪声时，分数可能很大，
实际算法还需要明确终止噪声、最后一步去噪与数据量化的处理。
